\section{Trabajo futuro}
\label{sec:futuro}

Se ordena por el criterio que gobierna todo el proyecto: cada incorporación
tiene que resolver un problema observado, no sonar avanzada.

\subsection{Lo que arreglaría limitaciones ya conocidas}

\begin{description}[leftmargin=0.6cm,style=nextline]
  \item[Contador de limitación compartido.] En cuanto el despliegue pase de un
    proceso, el límite actual deja de significar lo que dice. Es la condición
    escrita de reingreso de Redis, donde sí resolvería un problema real sin
    romper ninguna garantía transaccional.
  \item[Revocación de tokens.] Una lista de revocación consultada por
    identificador de token, o el paso a sesiones en base de datos si la
    revocación inmediata se vuelve un requisito.
  \item[Cookie \codigo{httpOnly} para el navegador.] Exige que la API y el
    frontend compartan origen, y protección contra falsificación de petición
    entre sitios.
  \item[Réplica de lectura.] Las consultas de saldo e historial no necesitan
    la instancia primaria. Antes de hacerlo habría que medir si la lectura es
    realmente el cuello de botella; con los datos actuales no lo es.
  \item[Medición de recursos.] Instrumentar CPU y memoria por proceso durante
    las carreras, para completar lo que la sección~\ref{sec:rendimiento} deja
    pendiente.
\end{description}

\subsection{Lo que ampliaría el alcance}

\begin{description}[leftmargin=0.6cm,style=nextline]
  \item[Reversiones contables.] Una transferencia errónea no se borra: se
    compensa con un asiento de reversión. El libro mayor de solo inserción ya
    está preparado para ello; falta el caso de uso y su autorización.
  \item[Varias cuentas por usuario.] El modelo lo admite y la autorización ya
    está escrita para el caso general; solo habría que levantar la
    restricción de unicidad y ajustar la interfaz.
  \item[Fracciones de céntimo.] Si alguna vez hicieran falta intereses o
    comisiones prorrateadas, la vía es \codigo{NUMERIC} con una política de
    redondeo explícita, nunca punto flotante.
  \item[Arquitectura orientada a eventos para lo posterior.] Notificaciones,
    puntuación de riesgo y analítica son consumidores naturales del libro
    mayor. La transferencia seguiría siendo síncrona y atómica.
\end{description}

\subsection{Inteligencia artificial, con condiciones}

El límite de diseño es permanente: \textbf{ningún modelo aprueba, rechaza ni
modifica una transferencia o un saldo}. Dentro de ese límite:

\begin{enumerate}[leftmargin=1.4em]
  \item \textbf{Reglas deterministas de riesgo} (velocidad de envío, importe
    atípico frente al historial propio, receptor nuevo) que adjunten una marca
    de revisión. Son baratas, explicables y auditables, y es lo que un sistema
    real pone primero.
  \item \textbf{Detección de anomalías no supervisada} como experimento, con
    la etiqueta explícita de que los datos son sintéticos y el resultado no es
    generalizable.
  \item Un clasificador supervisado de fraude \textbf{no} entra: necesitaría
    transacciones reales etiquetadas, que no existen y no se pueden fabricar
    sin falsear el resultado~\cite{sculley2015}.
\end{enumerate}
